\chapter{CAR Hilbert Carriers and Sign Involutions}
\label{app:car_hilbert}

\section{Introduction and the Two-Tiered Architecture}

The canonical anticommutation relations (CAR) govern the algebra of fermionic creation and annihilation operators. In physics textbooks, algebraic anticommutation relations and Hilbert-space adjoints are frequently introduced simultaneously. However, in constructive mathematics and formal verification, conflating these two aspects introduces subtle vulnerabilities:
\begin{enumerate}
    \item \textbf{Non-Unitary Similarities:} Let $c_i, c_i^\dagger$ satisfy the algebraic CAR on a vector space $V$. For any invertible linear transformation $S \in \operatorname{GL}(V)$, the similarity-transformed operators:
    \begin{equation}
    \tilde{c}_i \coloneqq S c_i S^{-1}, \qquad \tilde{c}_i^\dagger \coloneqq S c_i^\dagger S^{-1}
    \end{equation}
    satisfy all algebraic anticommutation relations identically:
    \[
    \{\tilde{c}_i, \tilde{c}_j\} = 0, \qquad \{\tilde{c}_i^\dagger, \tilde{c}_j^\dagger\} = 0, \qquad \{\tilde{c}_i, \tilde{c}_j^\dagger\} = \delta_{ij} I.
    \]
    However, unless $S$ is unitary, $\tilde{c}_i^\dagger \neq (\tilde{c}_i)^\dagger$. Thus, algebraic CAR does not imply adjoint compatibility.
    \item \textbf{Carrier Metric Distinctions:} The function carrier $X \to \mathbb{C}$ carries a natural product norm, but not an intrinsic Euclidean $\ell^2$ inner product. Conflating function spaces with Euclidean spaces risks installing incompatible metric instances.
\end{enumerate}

To resolve these issues, we establish a strict two-tiered algebraic hierarchy:
\begin{itemize}
    \item \textbf{Tier 1: Abstract Algebraic CAR,} defined over an arbitrary complex vector space $V$ without any metric or topological structure.
    \item \textbf{Tier 2: Hilbert-Space CAR,} defined over finite-dimensional complex Euclidean spaces, adding the true Hilbert adjoint condition $c_i^\dagger = (c_i)^\dagger$.
\end{itemize}

\section{Tier 1: Abstract Algebraic CAR Relations}

\begin{definition}[Algebraic CAR System]
Let $V$ be a complex vector space and let $\iota$ be a finite index set. An \emph{algebraic CAR system} on $(V, \iota)$ consists of two operator families $c, c^\dagger : \iota \to \operatorname{End}_{\mathbb{C}}(V)$ satisfying:
\begin{enumerate}
    \item $\{c_i, c_j\} = c_i c_j + c_j c_i = 0$ for all $i, j \in \iota$.
    \item $\{c_i^\dagger, c_j^\dagger\} = c_i^\dagger c_j^\dagger + c_j^\dagger c_i^\dagger = 0$ for all $i, j \in \iota$.
    \item $\{c_i, c_j^\dagger\} = c_i c_j^\dagger + c_j^\dagger c_i = \delta_{ij} I_V$ for all $i, j \in \iota$.
\end{enumerate}
The local number operator for mode $i \in \iota$ is defined algebraically by:
\begin{equation}
n_i \coloneqq c_i^\dagger c_i.
\end{equation}
\end{definition}

Without assuming any inner product, the algebraic relations alone imply fundamental commutation theorems.

\begin{theorem}[CAR Bilinear Commutator Identity]
\label{thm:app_bilinear_commutator}
Let $(c, c^\dagger)$ be an algebraic CAR system on $(V, \iota)$. For any four indices $a, b, c, d \in \iota$:
\begin{equation}
\left[ c_a^\dagger c_b, \; c_c^\dagger c_d \right] = \delta_{bc} c_a^\dagger c_d - \delta_{ad} c_c^\dagger c_b.
\end{equation}
\end{theorem}

\begin{proof}
We evaluate the commutator $[A B, C D]$ for operators satisfying $\{B, C\} = \delta_{bc} I$, $\{D, A\} = \delta_{ad} I$, $\{A, C\} = 0$, and $\{B, D\} = 0$.
By non-commutative ring algebra:
\begin{align*}
[c_a^\dagger c_b, \; c_c^\dagger c_d] &= c_a^\dagger c_b c_c^\dagger c_d - c_c^\dagger c_d c_a^\dagger c_b \\
&= c_a^\dagger (\delta_{bc} I - c_c^\dagger c_b) c_d - c_c^\dagger (\delta_{ad} I - c_a^\dagger c_d) c_b \\
&= \delta_{bc} c_a^\dagger c_d - \delta_{ad} c_c^\dagger c_b - c_a^\dagger c_c^\dagger c_b c_d + c_c^\dagger c_a^\dagger c_d c_b.
\end{align*}
Using $\{c_a^\dagger, c_c^\dagger\} = 0$ and $\{c_b, c_d\} = 0$:
\[
c_c^\dagger c_a^\dagger c_d c_b = (- c_a^\dagger c_c^\dagger) (- c_b c_d) = c_a^\dagger c_c^\dagger c_b c_d.
\]
The last two terms cancel identically:
\[
- c_a^\dagger c_c^\dagger c_b c_d + c_a^\dagger c_c^\dagger c_b c_d = 0.
\]
This leaves exactly $\delta_{bc} c_a^\dagger c_d - \delta_{ad} c_c^\dagger c_b$.
\end{proof}

\begin{corollary}[Number Operator Commutators and Idempotence]
\label{cor:app_number_comm}
For all modes $i, j \in \iota$:
\begin{enumerate}
    \item $[n_i, c_j^\dagger] = \delta_{ij} c_j^\dagger$.
    \item $[n_i, c_j] = - \delta_{ij} c_j$.
    \item $[n_i, n_j] = 0$.
    \item $n_i^2 = n_i$ (idempotence).
\end{enumerate}
\end{corollary}

\begin{proof}
\begin{enumerate}
    \item Set $a = i, b = i, c = j$ in the algebraic expansion:
    \[
    [c_i^\dagger c_i, c_j^\dagger] = c_i^\dagger \{c_i, c_j^\dagger\} - \{c_i^\dagger, c_j^\dagger\} c_i = c_i^\dagger (\delta_{ij} I) - 0 = \delta_{ij} c_j^\dagger.
    \]
    \item Similarly, $[c_i^\dagger c_i, c_j] = c_i^\dagger \{c_i, c_j\} - \{c_i^\dagger, c_j\} c_i = 0 - (\delta_{ij} I) c_i = - \delta_{ij} c_i$.
    \item Setting $a = b = i$ and $c = d = j$ in Theorem~\ref{thm:app_bilinear_commutator}:
    \[
    [n_i, n_j] = [c_i^\dagger c_i, c_j^\dagger c_j] = \delta_{ij} c_i^\dagger c_j - \delta_{ij} c_j^\dagger c_i = \delta_{ij} (n_i - n_i) = 0.
    \]
    \item Setting $i = j$ in the anticommutator $\{c_i, c_i^\dagger\} = I$:
    \[
    n_i^2 = (c_i^\dagger c_i)(c_i^\dagger c_i) = c_i^\dagger (I - c_i^\dagger c_i) c_i = c_i^\dagger c_i - (c_i^\dagger)^2 c_i^2.
    \]
    Since $2 (c_i^\dagger)^2 = \{c_i^\dagger, c_i^\dagger\} = 0$, we have $(c_i^\dagger)^2 = 0$, so $n_i^2 = c_i^\dagger c_i = n_i$.
\end{enumerate}
\end{proof}

\section{Tier 2: Concrete Occupation Hilbert Spaces}

We now construct the concrete representation of CAR on the fermionic Fock space.

\begin{definition}[Fermionic Fock Space]
Let $\iota$ be a finite index set of cardinality $|\iota| = d$. The \emph{fermionic Fock space} $\mathcal{F}(\iota)$ is the complex Euclidean space:
\begin{equation}
\mathcal{F}(\iota) \coloneqq \ell^2(\mathcal{P}(\iota)),
\end{equation}
where $\mathcal{P}(\iota)$ is the power set of all subsets of $\iota$.
The canonical orthonormal basis vectors are the occupation kets:
\begin{equation}
|S\rangle \in \mathcal{F}(\iota) \quad \text{for each } S \subseteq \iota.
\end{equation}
\end{definition}

\begin{proposition}[Carrier Properties]
\label{prop:fock_carrier_properties}
\begin{enumerate}
    \item \textbf{Finite Dimension:} $\dim_{\mathbb{C}} \mathcal{F}(\iota) = 2^{|\iota|}$.
    \item \textbf{Inner Product:} For $u, v \in \mathcal{F}(\iota)$, $\langle u \mid v \rangle = \sum_{S \subseteq \iota} \overline{u(S)} v(S)$.
    \item \textbf{Orthonormality:} $\langle S \mid T \rangle = \delta_{S, T}$ for all $S, T \subseteq \iota$.
    \item \textbf{Basis Resolution:} Every vector $v \in \mathcal{F}(\iota)$ expands uniquely as $v = \sum_{S \subseteq \iota} v(S) |S\rangle$.
\end{enumerate}
\end{proposition}

\begin{lemma}[Adjoint Verification via Basis Pairings]
\label{lem:adjoint_pairing}
Let $A, B \in \operatorname{End}_{\mathbb{C}}(\mathcal{F}(\iota))$ be linear operators. If the inner product pairing satisfies:
\begin{equation}
\langle S \mid A T \rangle = \langle B S \mid T \rangle \quad \text{for all basis pairs } S, T \subseteq \iota,
\end{equation}
then $B$ is identically the Hilbert adjoint of $A$:
\begin{equation}
B = A^\dagger.
\end{equation}
\end{lemma}

\begin{proof}
Let $u = \sum_S u(S) |S\rangle$ and $v = \sum_T v(T) |T\rangle$ be arbitrary vectors in $\mathcal{F}(\iota)$. By sesquilinearity:
\begin{align*}
\langle u \mid A v \rangle &= \sum_{S, T \subseteq \iota} \overline{u(S)} v(T) \langle S \mid A T \rangle \\
&= \sum_{S, T \subseteq \iota} \overline{u(S)} v(T) \langle B S \mid T \rangle = \langle B u \mid v \rangle.
\end{align*}
By the Riesz representation characterization of the adjoint, this proves $B = A^\dagger$.
\end{proof}

\section{Fermionic Sign Mechanics and Jordan--Wigner Ordering}

To construct the operators $c_i, c_i^\dagger$, we fix a strict linear order $<$ on $\iota$.

\begin{definition}[Preceding Occupation Count and Sign]
For any mode $i \in \iota$ and configuration $S \subseteq \iota$:
\begin{enumerate}
    \item The \emph{preceding occupied count} is the number of occupied modes strictly below $i$:
    \begin{equation}
    \sigma(i, S) \coloneqq \#\{ j \in S \mid j < i \}.
    \end{equation}
    \item The \emph{fermionic sign factor} is:
    \begin{equation}
    \epsilon(i, S) \coloneqq (-1)^{\sigma(i, S)} \in \{+1, -1\}.
    \end{equation}
\end{enumerate}
\end{definition}

To prevent natural number underflow in formal induction, the preceding count is governed by strictly additive recursion formulas.

\begin{lemma}[Additive Preceding Count Laws]
\label{lem:additive_count}
Let $i, j \in \iota$ and $S \subseteq \iota$.
\begin{enumerate}
    \item \textbf{Insertion Formula:} For $j \notin S$:
    \begin{equation}
    \sigma(i, S \cup \{j\}) = \sigma(i, S) + \begin{cases} 1, & \text{if } j < i, \\ 0, & \text{if } j \ge i. \end{cases}
    \end{equation}
    \item \textbf{Erasure Formula:} For $j \in S$:
    \begin{equation}
    \sigma(i, S) = \sigma(i, S \setminus \{j\}) + \begin{cases} 1, & \text{if } j < i, \\ 0, & \text{if } j \ge i. \end{cases}
    \end{equation}
    \item \textbf{Self Invariance:} $\sigma(i, S \cup \{i\}) = \sigma(i, S)$ and $\sigma(i, S \setminus \{i\}) = \sigma(i, S)$.
\end{enumerate}
\end{lemma}

\begin{proof}
The set $\{k \in S \cup \{j\} \mid k < i\}$ partitions into $\{k \in S \mid k < i\}$ and $\{j\}$ if $j < i$, or just $\{k \in S \mid k < i\}$ if $j \ge i$. Taking cardinalities yields the additive formula. The erasure formula follows by writing $S = (S \setminus \{j\}) \cup \{j\}$ and applying insertion.
\end{proof}

\begin{theorem}[Permutation Sign Exchange Identities]
\label{thm:app_sign_exchange}
For any two distinct modes $i \neq j$ in $\iota$:
\begin{enumerate}
    \item If $i, j \notin S$:
    \begin{equation}
    \epsilon(j, S) \epsilon(i, S \cup \{j\}) = - \epsilon(i, S) \epsilon(j, S \cup \{i\}).
    \end{equation}
    \item If $i, j \in S$:
    \begin{equation}
    \epsilon(j, S) \epsilon(i, S \setminus \{j\}) = - \epsilon(i, S) \epsilon(j, S \setminus \{i\}).
    \end{equation}
    \item If $i \notin S$ and $j \in S$:
    \begin{equation}
    \epsilon(j, S) \epsilon(i, S \setminus \{j\}) = - \epsilon(i, S) \epsilon(j, S \cup \{i\}).
    \end{equation}
\end{enumerate}
\end{theorem}

\begin{proof}
Without loss of generality, assume $j < i$.
\begin{enumerate}
    \item By Lemma~\ref{lem:additive_count}(1), since $j < i$, $\sigma(i, S \cup \{j\}) = \sigma(i, S) + 1$.
    On the other hand, since $i > j$, $\sigma(j, S \cup \{i\}) = \sigma(j, S) + 0 = \sigma(j, S)$.
    Therefore:
    \[
    \epsilon(j, S) \epsilon(i, S \cup \{j\}) = (-1)^{\sigma(j, S)} (-1)^{\sigma(i, S) + 1} = - (-1)^{\sigma(j, S) + \sigma(i, S)}.
    \]
    While the right-hand side has:
    \[
    \epsilon(i, S) \epsilon(j, S \cup \{i\}) = (-1)^{\sigma(i, S)} (-1)^{\sigma(j, S)} = (-1)^{\sigma(i, S) + \sigma(j, S)}.
    \]
    The ratio of the two products is $(-1)^1 = -1$.
    \item The proofs for items 2 and 3 follow identically from Lemma~\ref{lem:additive_count} by tracking the single relative shift $[j < i] = 1$.
\end{enumerate}
\end{proof}

\begin{definition}[Concrete Creation and Annihilation Operators]
The linear operators $c_i, c_i^\dagger \in \operatorname{End}_{\mathbb{C}}(\mathcal{F}(\iota))$ are defined on basis kets $|S\rangle$ by:
\begin{align}
c_i |S\rangle &\coloneqq \begin{cases} \epsilon(i, S) |S \setminus \{i\}\rangle, & \text{if } i \in S, \\ 0, & \text{if } i \notin S, \end{cases} \\
c_i^\dagger |S\rangle &\coloneqq \begin{cases} \epsilon(i, S) |S \cup \{i\}\rangle, & \text{if } i \notin S, \\ 0, & \text{if } i \in S. \end{cases}
\end{align}
\end{definition}

\begin{theorem}[Concrete CAR and Adjoint Verification]
\label{thm:app_concrete_car}
The operators $(c, c^\dagger)$ satisfy:
\begin{enumerate}
    \item $\{c_i, c_j\} = 0$ for all $i, j \in \iota$.
    \item $\{c_i^\dagger, c_j^\dagger\} = 0$ for all $i, j \in \iota$.
    \item $\{c_i, c_j^\dagger\} = \delta_{ij} I$ for all $i, j \in \iota$.
    \item $c_i^\dagger = (c_i)^\dagger$ for all $i \in \iota$.
\end{enumerate}
\end{theorem}

\begin{proof}
\begin{enumerate}
    \item For $i = j$: $c_i^2 |S\rangle = 0$ because $i \notin S \setminus \{i\}$.
    For $i \neq j$: if both $i, j \in S$, applying $c_i c_j$ yields $\epsilon(j, S) \epsilon(i, S \setminus \{j\}) |S \setminus \{i, j\}\rangle$.
    Applying $c_j c_i$ yields $\epsilon(i, S) \epsilon(j, S \setminus \{i\}) |S \setminus \{i, j\}\rangle$.
    By Theorem~\ref{thm:app_sign_exchange}(2), the two coefficients sum to zero. If either $i \notin S$ or $j \notin S$, both compositions vanish.
    \item Identical to item 1 using Theorem~\ref{thm:app_sign_exchange}(1).
    \item For $i \neq j$: follows from Theorem~\ref{thm:app_sign_exchange}(3).
    For $i = j$: if $i \in S$, $c_i^\dagger c_i |S\rangle = 1 |S\rangle$ and $c_i c_i^\dagger |S\rangle = 0$.
    If $i \notin S$, $c_i^\dagger c_i |S\rangle = 0$ and $c_i c_i^\dagger |S\rangle = 1 |S\rangle$.
    In both cases, $(c_i c_i^\dagger + c_i^\dagger c_i)|S\rangle = |S\rangle$.
    \item To verify adjointness, let $S, T \subseteq \iota$.
    $\langle S \mid c_i T \rangle$ is non-zero only when $i \in T$ and $S = T \setminus \{i\}$, in which case it equals $\epsilon(i, T)$.
    Conversely, $\langle c_i^\dagger S \mid T \rangle$ is non-zero only when $i \notin S$ and $S \cup \{i\} = T$, in which case it equals $\epsilon(i, S)$.
    By Lemma~\ref{lem:additive_count}(3), $\epsilon(i, T) = \epsilon(i, S \cup \{i\}) = \epsilon(i, S)$.
    Therefore $\langle S \mid c_i T \rangle = \langle c_i^\dagger S \mid T \rangle$.
    By Lemma~\ref{lem:adjoint_pairing}, $c_i^\dagger = (c_i)^\dagger$.
\end{enumerate}
\end{proof}

\section{Global Parity and Ordered Product Mechanics}

Because operator composition on $\operatorname{End}_{\mathbb{C}}(\mathcal{F}(\iota))$ is non-commutative, infinite or finite products of operators cannot be written without specifying a linear order.

\begin{definition}[Global Parity Operator]
Let $\operatorname{sort}(\le, \iota) = (i_1, i_2, \dots, i_d)$ be the sorted list of all modes in $\iota$. The \emph{global parity operator} $\Gamma$ is:
\begin{equation}
\Gamma \coloneqq \prod_{k=1}^d (I - 2 n_{i_k}).
\end{equation}
\end{definition}

\begin{theorem}[Parity Properties]
\label{thm:app_parity_properties}
\begin{enumerate}
    \item On basis kets: $\Gamma |S\rangle = (-1)^{|S|} |S\rangle$.
    \item Involution and Hermiticity: $\Gamma^2 = I$ and $\Gamma^\dagger = \Gamma$.
    \item Anticommutation with Fermions: $\{\Gamma, c_i\} = 0$ and $\{\Gamma, c_i^\dagger\} = 0$ for all $i \in \iota$.
\end{enumerate}
\end{theorem}

\begin{proof}
\begin{enumerate}
    \item Since $[n_i, n_j] = 0$, the product acts as the product of eigenvalues $(1 - 2 [i \in S])$. Each occupied mode contributes $(1 - 2) = -1$, while unoccupied modes contribute $(1 - 0) = 1$. The product is $(-1)^{|S|}$.
    \item $(-1)^{|S|} (-1)^{|S|} = 1$, so $\Gamma^2 = I$. Since eigenvalues are real, $\Gamma^\dagger = \Gamma$.
    \item $c_i$ maps $|S\rangle$ (parity $(-1)^{|S|}$) to $|S \setminus \{i\}\rangle$ (parity $(-1)^{|S|-1} = -(-1)^{|S|}$). Thus $\Gamma c_i = - c_i \Gamma$.
\end{enumerate}
\end{proof}

\begin{remark}[Formalization Note]
In the verified Lean~4 development (Appendix~A02), all sign exchange lemmas, the basis extension constructor, and the adjoint pairing theorems are verified constructively. In particular, the absence of natural underflow ensures that sign proofs do not require classical casing over integer signs.
\end{remark}
